% =====================================================================
% Supplementary Material
% =====================================================================

\section*{Supplementary Material}

This supplementary document provides additional results and analysis that complement the main paper. We organize the material as follows: Section~\ref{sec:sup_supervised} presents a reference comparison with supervised post-training methods, Section~\ref{sec:sup_expanded} reports expanded understanding benchmarks including specialist evaluations, Section~\ref{sec:sup_cost} analyzes computational overhead, Section~\ref{sec:sup_additional_qualitative} provides additional qualitative examples, and Section~\ref{sec:sup_per_model_dynamics} presents per-model training dynamics.


% -------------------------------------------------------------------
% S1: Supervised Reference Comparison
% -------------------------------------------------------------------
\subsection*{S1. Reference Comparison with Supervised Methods}
\label{sec:sup_supervised}

Table~\ref{tab:supervised_reference} compares our framework against post-training methods that employ curated fine-tuning data, external reward models, or supervised initialization stages. These methods are not directly comparable to our fully unsupervised approach, as they benefit from additional human supervision. We include them as reference upper bounds to contextualize the gap between supervised and self-supervised post-training and to demonstrate that unsupervised self-evolution can approach supervised performance levels without any labeled data.

\input{tables/supplementary_tables}


% -------------------------------------------------------------------
% S2: Expanded Understanding Benchmarks
% -------------------------------------------------------------------
\subsection*{S2. Expanded Understanding Benchmarks}
\label{sec:sup_expanded}

Table~\ref{tab:expanded_understanding} extends the main paper's understanding evaluation to include OCRBench (optical character recognition) and MathVista (mathematical visual reasoning). These specialist benchmarks test capabilities outside the training distribution, as our framework uses only general visual questions and does not specifically target OCR or mathematical reasoning. Maintaining or improving performance on these benchmarks provides evidence that the self-evolving loop does not degrade specialist capabilities while improving general understanding.


% -------------------------------------------------------------------
% S3: Computational Cost Analysis
% -------------------------------------------------------------------
\subsection*{S3. Computational Cost Analysis}
\label{sec:sup_cost}

% TODO: Fill with actual timing and resource numbers once training is complete.

\begin{table}[tb]
  \caption{Computational overhead of the self-evolving framework. Wall-clock times are measured on NVIDIA A100 GPUs. STE adds negligible cost (single greedy forward pass) relative to multi-sample self-consistency evaluation.}
  \label{tab:computational_cost}
  \centering
  \small
  \begin{tabular}{lcccc}
    \toprule
    Component & Forward Passes & Backward Passes & Time per Step & \% of Total \\
    \midrule
    Proposer (question generation) & -- & -- & -- & -- \\
    Prompt-perturbed self-consistency ($N{=}7$) & -- & -- & -- & -- \\
    STE computation (single greedy pass) & -- & -- & -- & -- \\
    Proposer GRPO update & -- & -- & -- & -- \\
    Solver REINFORCE update & -- & -- & -- & -- \\
    \midrule
    Generator (image synthesis) & -- & -- & -- & -- \\
    QA fidelity scoring ($M$ questions $\times$ 2 images) & -- & -- & -- & -- \\
    Cycle-consistent captioning & -- & -- & -- & -- \\
    Generator REINFORCE update & -- & -- & -- & -- \\
    \midrule
    Total per training cycle ($U{+}G$ steps) & -- & -- & -- & -- \\
    \bottomrule
  \end{tabular}
\end{table}

% Key points to cover:
% - STE is essentially free (1 forward pass vs N=7 for PPS)
% - Generation loop is more expensive than understanding loop due to image synthesis
% - Total training time compared to base model pretraining is small (post-training)
% - Memory overhead of 3 LoRA adapters is minimal


% -------------------------------------------------------------------
% S4: Additional Qualitative Examples
% -------------------------------------------------------------------
\subsection*{S4. Additional Qualitative Examples}
\label{sec:sup_additional_qualitative}

% TODO: Add additional qualitative comparison figure with more diverse examples.

\begin{figure*}[tb]
  \centering
  % \includegraphics[width=\linewidth]{figures/supplementary_qualitative.pdf}
  \fbox{\parbox{0.95\linewidth}{\centering\vspace{4cm}\textbf{[Additional Qualitative Examples Placeholder]}\\\smallskip
  Extended generation comparison across diverse prompts and all three architectures.\\
  Include failure cases with analysis.\vspace{4cm}}}
  \caption{Extended qualitative comparison across all three architectures. Each row shows the text prompt, the base model generation, and the self-evolved model generation. The final row presents failure cases where the self-evolving framework did not produce meaningful improvement, along with analysis of the underlying cause.}
  \label{fig:sup_qualitative}
\end{figure*}


% -------------------------------------------------------------------
% S5: Per-Model Training Dynamics
% -------------------------------------------------------------------
\subsection*{S5. Per-Model Training Dynamics}
\label{sec:sup_per_model_dynamics}

% TODO: Add per-model training dynamics figure.

\begin{figure*}[tb]
  \centering
  % \includegraphics[width=\linewidth]{figures/per_model_dynamics.pdf}
  \fbox{\parbox{0.95\linewidth}{\centering\vspace{3.5cm}\textbf{[Per-Model Training Dynamics Placeholder]}\\\smallskip
  Three columns (BLIP3o-8B, BAGEL, VARGPT) showing STE difficulty and solver accuracy curves.\\\smallskip
  Demonstrates that the self-evolving loop operates consistently across architectures.\vspace{3.5cm}}}
  \caption{Training dynamics for each backbone architecture. Despite fundamentally different generation mechanisms, all three models exhibit similar training dynamics: STE difficulty increases monotonically, bucket proportions shift from easy-dominated to balanced distributions, and solver accuracy improves steadily. The consistency of these dynamics across architectures provides additional evidence for the model-agnostic nature of the framework.}
  \label{fig:sup_per_model_dynamics}
\end{figure*}

% Key points:
% - Similar training dynamics across all three architectures
% - Convergence rates may differ but qualitative behavior is consistent
% - This supports the model-agnostic claim in the main paper
